\section{Hash functions}\label{sec:prelims-hash}

Hash functions compress inputs of arbitrary length into a short digest.

\paragraph{Syntax.} A hash function family consists of algorithms defined as
follows:
%
\begin{description}

  \item[$\key \sample \hashgen(\secparam)$ :] The key generation algorithm
    takes the security parameter as input, and outputs a key $\key$.

  \item[$y \leftarrow \hasheval(\key, x)$ :] The evaluation algorithm takes key
    $\key$ and input $x \in \{0,1\}^*$, and outputs digest $y$ in the range
    $\range$.

\end{description}
%
We say $\Hash = (\hashgen, \hash)$ is a hash function family with range
$\range$ if $\hasheval(\key, \cdot)$ is efficiently computable for every $\key$
output by $\hashgen(\secparam)$. We use the standard syntax and security
notions for hash functions~\cite{KatLin14}. When the key $\key$ is fixed and
public, as is the case throughout this thesis, we write $\hash(x)$ for
$\hasheval(\key, x)$.

\paragraph{Collision resistance.} Collision resistance requires that no
efficient adversary can find two distinct inputs hashing to the same digest.

\begin{experimentbox}{Collision resistance}{hash-cr}

  \begin{enumerate}

    \item The challenger generates $\key \sample \hashgen(\secparam)$ and sends
      $\key$ to $\adv$.

    \item The adversary outputs $(x, x')$.

  \end{enumerate}

  $\adv$ wins the experiment if $x \neq x'$ and $\hasheval(\key, x) =
  \hasheval(\key, x')$.

\end{experimentbox}

The collision resistance advantage of $\adv$ is defined as follows:
%
\begin{equation*}
%
  \Adv^{\mathsf{cr}}_{\Hash, \adv}(\secparam) = \Pr[\adv \text{ wins}].
%
\end{equation*}
%
We say $\Hash$ is collision resistant if this advantage is small for every
admissible adversary $\adv$.

\paragraph{Preimage resistance.} Preimage resistance requires that no efficient
adversary can recover an input hashing to a given digest, even one of their own
choosing.

\begin{experimentbox}{Preimage resistance}{hash-preimage}

  \begin{enumerate}

    \item The challenger generates $\key \sample \hashgen(\secparam)$, samples
      $x \sample \{0,1\}^{\secparam}$, computes $y \leftarrow \hasheval(\key,
      x)$, and sends $(\key, y)$ to $\adv$.

    \item The adversary outputs $x'$.

  \end{enumerate}

  $\adv$ wins the experiment if $\hasheval(\key, x') = y$.

\end{experimentbox}

The preimage resistance advantage of $\adv$ is defined as follows:
%
\begin{equation*}
%
  \Adv^{\mathsf{preimage}}_{\Hash, \adv}(\secparam) = \Pr[\adv \text{ wins}].
%
\end{equation*}
%
We say $\Hash$ is preimage resistant if this advantage is small for every
admissible adversary $\adv$. A related property is second preimage resistance,
which requires that no efficient adversary given $x$, can find a distinct $x'$
with $\hasheval(\key, x) = \hasheval(\key, x')$. Collision resistance implies
second preimage resistance, so we do not use it separately.

\paragraph{The random oracle model.} Our security proofs generally model a hash
function as a random oracle: a function sampled uniformly at random from all
functions with the given domain and range, to which every party, including the
adversary, has oracle access. The random oracle answers each fresh query with
an output of its choosing. It has internal state and records the message and
output for each query it has answered, which allows answering consistently for
repeated queries. We refer to the ability to choose the answer to oracle
queries as programming the random oracle.
